%% ---- sezione C: modelli+dati ----
\section{Dalla PCA al VAE: i modelli}

Il capitolo 4 presenta, in ordine di complessità crescente, i quattro modelli confrontati.

\paragraph{PCA.} È il riferimento lineare: una proiezione ortogonale $\mathbf{z} = W^{\top} \mathbf{x}$, $\hat{\mathbf{x}} = W W^{\top} \mathbf{x}$, ottima in termini di errore quadratico medio, con soluzione analitica negli autovettori dominanti della covarianza. Il limite è strutturale: assume che la manifold dei dati sia un iperpiano lineare, mentre i co-movimenti di mercato (tail dependence, gerarchie settoriali) presentano strutture non lineari.

\paragraph{Autoencoder lineare.} Con attivazioni identità, senza bias né weight decay, l'autoencoder a un solo strato converge, per il risultato di Bourlard e Kamp (1988), al sottospazio delle prime $K$ componenti principali: serve da sanity check e benchmark strutturale di partenza.

\paragraph{Autoencoder profondo.} Strati nascosti non lineari possono catturare la curvatura della manifold, ma lo spazio latente, non vincolato ad alcun prior e organizzato dal solo errore di ricostruzione, risulta irregolare, con ampie zone vuote. Il limite non sta nella qualità della compressione ma nell'obiettivo di addestramento, che non prevede alcun campionamento di nuovi punti nello spazio latente, operazione per la quale il modello non è addestrato (cap.~7); si tratta di una premessa metodologica, non sottoposta a verifica empirica nella tesi.

\paragraph{VAE.} Il Variational Autoencoder (Kingma e Welling, 2013) fa emettere all'encoder una distribuzione $q_\phi(\mathbf{z}|\mathbf{x}) = \mathcal{N}(\boldsymbol{\mu}, \mathrm{diag}(\boldsymbol{\sigma}^2))$, campionata con il reparameterization trick $\mathbf{z} = \boldsymbol{\mu} + \boldsymbol{\sigma} \odot \boldsymbol{\epsilon}$, $\boldsymbol{\epsilon} \sim \mathcal{N}(0, I)$. La loss somma un termine di ricostruzione, che nel cap.~6 è la divergenza di Jeffreys fra matrici di correlazione, e il regolarizzatore $\beta\, D_{\mathrm{KL}}(q_\phi(\mathbf{z}|\mathbf{x})\,\|\,\mathcal{N}(0, I))$ nello spazio latente, con $\beta = 1$. La KL addensa le distribuzioni attorno all'origine e le sovrappone, riducendo le zone vuote dello spazio latente. Rispetto alle GAN di CorrGAN (Marti, 2020), il VAE offre un addestramento tipicamente più stabile, privo del gioco minimax, e uno spazio latente esplicitamente strutturato e ricampionabile, proprietà sfruttata anche da Caprioli et al. (2024) per costruire scenari di rischio interpretabili. Il campionamento dal prior, possibile in linea di principio, è sostituito nella tesi dal block-bootstrap di incrementi storici nello spazio latente (cap.~8), ancorato allo stato di mercato osservato.

\section{Dati e pre-processing}

\paragraph{Dataset.} Il dataset raccoglie i prezzi di chiusura rettificati dei titoli dell'S\&P 500 dal 2000 al 2020, periodo che include fasi espansive, recessive e di stress. Sono mantenuti i soli titoli con quotazioni ininterrotte sull'intero arco, $N = 362$: la tesi dichiara il survivorship bias e legge il campione come ``popolazione condizionata alla sopravvivenza''. La classificazione GICS è stata ricostruita con un modello linguistico (gemini-2.5-flash, few-shot): 360 titoli su 362, con ABC e CERN non classificati.

\paragraph{Finestre mobili.} La correlazione è stimata sui log-rendimenti, stazionari a differenza dei prezzi, su una finestra mobile di $T_w = 724$ giorni con passo $S = 10$. Il valore $T_w = 2N$ fissa per costruzione il rapporto di spettro $q = N/T_w = 0.5$, il regime di Marchenko-Pastur del cap.~3, e assicura che ogni matrice empirica sia di rango pieno, cioè definita positiva, condizione di esistenza e unicità della decomposizione di Cholesky. Dai 5333 log-rendimenti si ottengono 461 matrici.

\paragraph{Filtro di Perron-Frobenius.} Le prime 49 finestre violano l'uniformità di segno del primo autovettore. Sono tutte e sole quelle con osservazioni antecedenti la metà di ottobre 2001 e il solo titolo con componente negativa è sempre NEM (Newmont), unico titolo aurifero del paniere, che durante la bolla dot-com si è mosso da bene rifugio, in controtendenza rispetto al listino. La violazione è dunque riconducibile a un regime eccezionale e circoscritto, non a un artefatto di misura. L'esclusione è dichiarata come scelta di perimetro: il lavoro si limita al regime cooperativo ``standard'', su cui sono definiti i fatti stilizzati del protocollo di validazione, senza pretendere di rappresentare i regimi di disaccoppiamento. Restano 412 matrici.

\paragraph{Ripartizione.} Test preliminari con split cronologico hanno evidenziato un distribution shift: i regimi passati non bastavano a generalizzare sui successivi. La tesi adotta perciò un random shuffle dell'intero bacino, ripartito 70/20/10 in 288, 82 e 42 matrici, così da distribuire omogeneamente fasi volatili e stabili nei tre set. Lo shuffle non elimina però la sovrapposizione fra finestre adiacenti, fino al 98.6\% delle osservazioni (714 giorni su 724): matrici quasi identiche possono finire in set diversi, limite dichiarato.

\paragraph{Parametrizzazione di Cholesky.} Un primo approccio usava il triangolo inferiore stretto della matrice di correlazione ($D = 65\,341$), con clipping in $[-1, 1]$. È stato abbandonato perché non offre alcuna garanzia di semidefinita positività, proprietà globale dello spettro, non imponibile coefficiente per coefficiente: le verifiche spettrali hanno rilevato autovalori negativi, e uno solo basta a far fallire la decomposizione di Cholesky. La parametrizzazione definitiva vettorizza il fattore $L$ di $C = L L^{\top}$, diagonale inclusa, con $D = N(N+1)/2 = 65\,703$. Per qualunque output del decoder, $\hat{L}\hat{L}^{\top}$ è una matrice di Gram, dunque semidefinita positiva per costruzione, a prescindere dalla qualità dell'apprendimento; la rinormalizzazione a diagonale unitaria $\hat{R} = \Delta^{-1/2}\hat{C}\Delta^{-1/2}$, con $\Delta = \mathrm{diag}(\hat{C})$, preserva la semidefinita positività per la legge d'inerzia di Sylvester. Lo stesso vettore $\mathbf{x} \in \mathbb{R}^{65\,703}$ è l'input di tutti e quattro i modelli.

\begin{riquadro}{La pipeline in numeri}
5334 quotazioni giornaliere $\rightarrow$ 5333 log-rendimenti $\rightarrow$ 461 finestre (724 giorni, passo 10) $\rightarrow$ filtro di Perron-Frobenius: $-49$ $\rightarrow$ 412 matrici $\rightarrow$ shuffle 70/20/10 $\rightarrow$ 288 / 82 / 42 $\rightarrow$ fattore di Cholesky vettorizzato, $D = 65\,703$.
\end{riquadro}
